\section{组合并行：3D/4D 并行实践}

\subsection{并行策略全景对比}

\begin{figure}[H]
\centering
\includegraphics[width=0.95\textwidth]{slides-images/slide-052.jpg}
\caption{所有并行策略的特性总结表\protect\footnotemark}
\end{figure}
\footnotetext{Slides 第53页。}

\begin{table}[H]
\centering
\small
\begin{tabular}{lcccc}
\toprule
\textbf{策略} & \textbf{通信开销} & \textbf{内存扩展} & \textbf{消耗批次} & \textbf{适用链路} \\
\midrule
DDP + ZeRO-1 & 每批 $2|\theta|$ & 无 & 是 & 慢速也可 \\
FSDP (ZeRO-3) & 每层同步 & 线性 & 是 & 中速 \\
流水线并行 & 小(点对点) & 线性 & 是 & 慢速 \\
张量并行 & 大(每层 All-Reduce) & 线性 & 否 & 高速必需 \\
\bottomrule
\end{tabular}
\caption{四种主要并行策略的关键特性对比}
\end{table}

\begin{importantbox}{三种核心有限资源}
大规模训练需要平衡三种有限资源：
\begin{enumerate}
    \item \textbf{内存}：参数 + 优化器状态 + 激活
    \item \textbf{带宽/计算}：通信开销 vs. 有效计算
    \item \textbf{批次大小}：一种非传统但极其重要的资源——数据并行和流水线并行都"消耗"它
\end{enumerate}
\end{importantbox}

\subsection{批次大小与并行策略的关系}

\begin{figure}[H]
\centering
\includegraphics[width=0.95\textwidth]{slides-images/slide-053.jpg}
\caption{Google TPU 团队的分析：不同批次大小下的最优并行组合\protect\footnotemark}
\end{figure}
\footnotetext{Slides 第54页。来自 Google TPU 并行化指南。}

\begin{knowledgebox}{批次大小决定最优策略}
\begin{itemize}
    \item \textbf{批次太小}：通信受限，无论用哪种策略都效率低下
    \item \textbf{中等批次}：需要 FSDP + 张量并行的组合才能达到计算受限（compute-bound）
    \item \textbf{大批次}：纯 FSDP（数据并行）就足以实现高效率
\end{itemize}
\end{knowledgebox}

\subsection{3D 并行的经验法则}

\begin{figure}[H]
\centering
\includegraphics[width=0.95\textwidth]{slides-images/slide-054.jpg}
\caption{3D 并行的设计经验法则\protect\footnotemark}
\end{figure}
\footnotetext{Slides 第55页。}

\begin{importantbox}{3D 并行配置法则}
按优先级从高到低：
\begin{enumerate}
    \item \textbf{首先用张量并行}填满节点内 GPU（通常 $T=8$），确保模型参数能放进内存
    \item \textbf{如果还放不下}，用 ZeRO Stage 3 或流水线并行跨节点扩展
    \item \textbf{剩余 GPU} 用数据并行来扩展总吞吐量
    \item \textbf{如果批次大小仍有余量}，用梯度累积增大有效批次，减少同步频率
\end{enumerate}
带宽需求从高到低：张量并行 $>$ 上下文并行 $>$ 流水线并行 $>$ 数据并行。
\end{importantbox}

\subsection{本章小结}

实际训练中需要组合多种并行策略。核心原则是：高通信量的策略放在高带宽链路上（张量并行 $\rightarrow$ 节点内），低通信量的策略放在慢速链路上（数据并行 $\rightarrow$ 跨机架）。批次大小是连接所有策略的关键资源。

%% ========== Part 9: 案例研究 ==========

